%===============================================================================
\section{Tổng quan các công trình liên quan và Phân tích khoảng trống}
%===============================================================================

\subsection{Khung kiến trúc MLOps và Mô hình trưởng thành}
Nghiên cứu của Amou Najafabadi \textit{et al.}~\cite{najafabadi2026} là công trình systematic mapping toàn diện nhất hiện nay, phân tích 93 bài báo khoa học để thiết lập khung kiến trúc MLOps chuẩn mực gồm 39 thành phần, 56 quy trình và 8 biến thể kiến trúc. Kreuzberger \textit{et al.}~\cite{kreuzberger2023} tổng hợp các nguyên lý vận hành, định nghĩa và phân tầng hạ tầng MLOps trong môi trường doanh nghiệp. Mô hình trưởng thành của Google Cloud~\cite{google_mlops} định nghĩa 3 cấp độ (ML0--ML2), trong đó Continuous Training nằm ở mức trưởng thành cao nhất ML2. Mô hình của Microsoft~\cite{microsoft_mlops} và nghiên cứu của Granlund \textit{et al.}~\cite{granlund2021} nhấn mạnh thách thức trong quản trị vòng đời và phân định trách nhiệm đa tổ chức. Tuy nhiên, các khung kiến trúc này dừng lại ở việc mô tả cấu trúc tĩnh hoặc quy định các trigger cơ học, chưa cung cấp giải pháp cụ thể về chẩn đoán nguyên nhân suy giảm và cơ chế ra quyết định thích ứng.

\subsection{Giám sát mô hình và Phát hiện trôi dạt}
Về mặt thuật toán học thuật, các cơ chế phát hiện trôi dạt kinh điển trên luồng dữ liệu (data streams) gồm ADWIN~\cite{bifet2007}, DDM~\cite{gama2004} và tổng quan toàn diện về concept drift của Lu \textit{et al.}~\cite{lu2019}. Ở khía cạnh công cụ công nghiệp, Evidently AI~\cite{evidently} cung cấp khả năng tính toán data drift và kiểm tra chất lượng dữ liệu cơ bản; NannyML~\cite{nannyml} tập trung vào phương pháp ước lượng hiệu năng (performance estimation) khi chưa có nhãn. Tuy vậy, các công cụ này hoạt động như các module đo lường độc lập, thiếu một vòng lặp khép kín có thể liên kết trực tiếp bằng chứng giám sát với các chính sách can thiệp an toàn và xử lý nhãn trễ thông qua mã định danh dự đoán (\codepath{prediction\_id}).

\subsection{Hệ thống Continuous Training và Tự thích ứng}
Luo \textit{et al.}~\cite{luo2023} đề xuất pipeline thử nghiệm thích ứng tự động (CTCD-e) dựa trên các công cụ nguồn mở nhằm tự động hóa quy trình huấn luyện và triển khai khi phát hiện drift. Semola \textit{et al.}~\cite{semola2023} giới thiệu Continual Learning as a Service (CLaaS) hỗ trợ cập nhật mô hình gia tăng trên điện toán đám mây. Đa số các hệ thống CT hiện nay đều thiết kế theo cơ chế phản xạ: \emph{khi có cảnh báo drift thì tự động kích hoạt huấn luyện lại}. Cách tiếp cận này bộc lộ hạn chế lớn vì chưa phân biệt được trôi dạt lành tính (Benign Drift) với trôi dạt có hại, dẫn đến lãng phí tài nguyên tính toán hoặc tái huấn luyện sai lầm trên dữ liệu rác.

\subsection{Quản trị rủi ro, Tính vững chắc và Vòng lặp chuyên gia (HITL)}
Bayram \& Ahmed~\cite{bayram2025} khảo sát toàn diện về tính vững chắc (robustness) và quản trị rủi ro trong MLOps, khẳng định việc tự động hóa hoàn toàn mà thiếu kiểm soát an toàn có thể gây thiệt hại nghiêm trọng cho doanh nghiệp. Các nghiên cứu gần đây bắt đầu nhấn mạnh vai trò của Human-in-the-Loop (HITL) trong việc kiểm chứng dữ liệu bất thường và giải quyết các trường hợp suy giảm hiệu năng không rõ nguyên nhân.

\subsection{Bảng đánh giá khoảng trống nghiên cứu (Literature Gap Matrix)}

Bảng~\ref{tab:lit-gap} đối chiếu toàn diện hiện trạng y văn và nền tảng hiện hữu với mục tiêu thiết kế của đề tài nghiên cứu theo 7 tiêu chí chức năng cốt lõi:

\begin{xltabular}{\linewidth}{@{}L{2.8cm} C{0.7cm} C{0.7cm} C{0.7cm} C{0.7cm} C{0.7cm} C{0.7cm} C{0.7cm} Y@{}}
\caption{Bảng đối chiếu và đánh giá khoảng trống nghiên cứu.}\label{tab:lit-gap}\\
\toprule
\textbf{Công trình / Nền tảng} & \textbf{C1} & \textbf{C2} & \textbf{C3} & \textbf{C4} & \textbf{C5} & \textbf{C6} & \textbf{C7} & \textbf{Khoảng trống chưa giải quyết (Identified Research Gap)} \\
\midrule
\endfirsthead
\toprule
\textbf{Công trình / Nền tảng} & \textbf{C1} & \textbf{C2} & \textbf{C3} & \textbf{C4} & \textbf{C5} & \textbf{C6} & \textbf{C7} & \textbf{Khoảng trống chưa giải quyết (Identified Research Gap)} \\
\midrule
\endhead

Amou Najafabadi \textit{et al.}~\cite{najafabadi2026} & \pmark & \pmark & \xmark & \xmark & \xmark & \xmark & \xmark & Khung tham chiếu tổng hợp mang tính mô tả; chưa hiện thực hóa thuật toán chẩn đoán hay chính sách CT thích ứng. \\
\addlinespace

Kreuzberger \textit{et al.}~\cite{kreuzberger2023} & \cmark & \pmark & \xmark & \xmark & \xmark & \pmark & \xmark & Định nghĩa CT ở cấp độ kiến trúc tổng quan; kích hoạt CT dạng phản xạ cơ học, thiếu chẩn đoán nguyên nhân và phân tầng Severity. \\
\addlinespace

Luo \textit{et al.} (CTCD-e)~\cite{luo2023} & \cmark & \pmark & \xmark & \xmark & \xmark & \pmark & \pmark & Tái huấn luyện tự động theo lịch/drift đơn lẻ; thiếu cơ chế phân loại trôi dạt lành tính vs. có hại và thiếu vòng lặp HITL. \\
\addlinespace

Semola \textit{et al.} (CLaaS)~\cite{semola2023} & \pmark & \pmark & \xmark & \xmark & \xmark & \pmark & \xmark & Tập trung thuật toán học liên tục (continual learning); giả định nhãn luôn sẵn sàng, thiếu chẩn đoán sự cố dữ liệu. \\
\addlinespace

Evidently AI / NannyML~\cite{evidently,nannyml} & \cmark & \cmark & \xmark & \xmark & \xmark & \xmark & \xmark & Công cụ đo lường và ước lượng trôi dạt độc lập; không có vòng lặp ra quyết định hành động hay điều phối vòng đời CT. \\
\addlinespace

Bayram \& Ahmed~\cite{bayram2025} & \pmark & \pmark & \xmark & \pmark & \pmark & \xmark & \pmark & Khảo sát lý thuyết về tính vững chắc và rủi ro; chưa đề xuất một framework vận hành và bộ benchmark định lượng cụ thể. \\
\addlinespace

ADWIN / DDM / Lu~\cite{bifet2007,gama2004,lu2019} & \cmark & \xmark & \xmark & \xmark & \xmark & \xmark & \xmark & Thuật toán phát hiện biến động trên luồng dữ liệu thô; hoàn toàn nằm ngoài bối cảnh kiến trúc MLOps và pipeline serving/CT. \\
\midrule

\textbf{Nền tảng MLOps hiện hữu} & \pmark & \pmark & \xmark & \xmark & \xmark & \pmark & \pmark & Đã có pipeline serving, tracking, model registry và kích hoạt huấn luyện cơ bản; thiếu chẩn đoán đa bằng chứng và chính sách thích ứng. \\
\addlinespace

\textbf{Mục tiêu thiết kế Đề xuất} & \cmark & \cmark & \cmark & \cmark & \cmark & \cmark & \cmark & \textbf{Khép kín chuỗi xử lý khoa học:} Giám sát đa nguồn $\rightarrow$ Chẩn đoán đa bằng chứng $\rightarrow$ Phân tầng Severity $\rightarrow$ Lựa chọn chiến lược thích ứng $\rightarrow$ Cổng thẩm định \& Triển khai an toàn. \\
\bottomrule
\end{xltabular}

\textit{Quy ước ký hiệu:} \cmark: Hỗ trợ đầy đủ / Mục tiêu thiết kế; \pmark: Hỗ trợ một phần / Mức khái niệm; \xmark: Chưa hỗ trợ / Khoảng trống nghiên cứu.

\textit{Ghi chú các tiêu chí đánh giá (C1--C7):}
\begin{itemize}
  \item \textbf{C1 (Multi-Signal Drift Detection):} Phát hiện đa tín hiệu (Data drift $P(X)$, Prediction drift $P(\hat{Y})$, Data quality).
  \item \textbf{C2 (Performance Monitoring):} Giám sát hiệu năng thực tế ($P(Y|X)$, F1, Precision, Recall).
  \item \textbf{C3 (Delayed-Label Handling):} Xử lý nhãn trễ dựa trên \codepath{prediction\_id} và cửa sổ trượt watermark.
  \item \textbf{C4 (Root-Cause Diagnosis):} Phân loại nguyên nhân suy giảm cốt lõi và định lượng mức độ nghiêm trọng ($S_0$--$S_4$).
  \item \textbf{C5 (Adaptive Action Space):} Không gian hành động đa nhánh (No-op, Incremental CT, Full CT, HITL).
  \item \textbf{C6 (Safe Evaluation Gate):} Cổng thẩm định đối kháng Champion--Candidate đa tiêu chí trước khi triển khai.
  \item \textbf{C7 (Traffic Experiment \& Rollback):} Hỗ trợ Shadow Mode, Canary rollout và cơ chế thu hồi tự động khi vi phạm rào chắn.
\end{itemize}

\subsection{Định vị Công bố Khoa học (Publication Positioning)}
Kết quả nghiên cứu của đề tài hướng tới đóng gói bài báo khoa học gửi đến các diễn đàn công bố quốc tế uy tín theo các tầng đóng góp:
\begin{itemize}
  \item \textbf{Đóng góp Cốt lõi (Core Contribution):} Các tạp chí và hội nghị uy tín chuyên sâu về hệ thống và phần mềm: Journal of Systems and Software (JSS), Automated Software Engineering (ASE), Empirical Software Engineering (EMSE), ICSME.
  \item \textbf{Mở rộng Chuyên sâu (Advanced Extension):} Các hội nghị và tạp chí chuyên ngành hàng đầu về thực nghiệm và tiến hóa phần mềm: EMSE, JSS, MSR (Mining Software Repositories).
  \item \textbf{Đột phá Dài hạn (Stretch / Future Direction):} Các diễn đàn hàng đầu thế giới về Công nghệ Phần mềm \& AI: IEEE Transactions on Software Engineering (TSE), ACM Transactions on Software Engineering and Methodology (TOSEM), ICSE, ESEC/FSE.
\end{itemize}
